% !PRINTED 72
\weberpage{72}

Aus (2) ergibt sich noch, mit Benutzung von (1) § 20, wenn
man $g_1' = 2$, $g_2' = 0$, oder $g_1' = 0$, $g_2' = 2$ setzt:

\weq{3}{
\begin{aligned}
\vartheta_{g_1+2,\,g_2}(v) &= e^{2\pi i g_2}\,\vartheta_{g_1,g_2}(v),\\
\vartheta_{g_1,\,g_2+2}(v) &= e^{\pi i g_1}\,\vartheta_{g_1,g_2}(v),
\end{aligned}
}

und durch (2) ist also zugleich die Periodizität der Funktionen
$\vartheta_{g_1,g_2}(v)$ vollständig ausgedrückt. Beispielsweise ergibt sich, wenn
$\mu$ und $\nu$ ganze Zahlen sind:

\weq{4}{
\vartheta_{00}\!\left(v - \frac{(2\nu-1)\,\omega + 2\mu - 1}{2}\right)
= (-1)^{\nu+1}\, i\, e^{\frac{-\pi i \omega}{4}(2\nu-1)^2}\, e^{\pi i (2\nu-1) v}\, \vartheta_{11}(v).
}

\weq{5}{
\vartheta_{11}(v - \nu\omega - \mu) = (-1)^{\mu+\nu}\, e^{-\pi i \omega \nu^2}\, e^{2\pi i \nu v}\, \vartheta_{11}(v).
}

Die Formel (2) ist gültig für ganz beliebige, selbst komplexe
Werte von $g_1, g_2$.

Wenn man in den Definitionsformeln der $\vartheta$-Funktionen
[§ 20, (1)] $v$ durch $-v-1$, und durch $-v-\omega$, ferner $g_1, g_2$ durch
$-g_1, -g_2$ ersetzt, so zeigt sich, daß $\vartheta_{-g_1,-g_2}(-v)$ denselben Bedingungen
 genügt, wie $\vartheta_{g_1,g_2}(v)$, und daß sonach

\[
\vartheta_{g_1,g_2}(v),\quad \vartheta_{-g_1,-g_2}(-v)
\]

bis auf einen konstanten Faktor identisch sind; es ist also insbesondere, wie aus $v = 0$ hervorgeht:

\[
\vartheta_{0,0}(v) = \vartheta_{0,0}(-v),
\]

und danach ergibt die Formel (1), wenn man $v, g_1, g_2$ durch
$-v, -g_1, -g_2$ ersetzt:

\weq{6}{
\vartheta_{g_1,g_2}(v) = \vartheta_{-g_1,-g_2}(-v).
}

Sind die Elemente $g_1, g_2$ ganze Zahlen, so heißt $(g_1, g_2)$
eine \gesperrt{Hauptcharakteristik}. Es gibt deren vier wesentlich
verschiedene, nämlich:

\[
(0,0),\quad (0,1),\quad (1,0),\quad (1,1),
\]

und demnach auch vier wesentlich verschiedene Haupt-$\vartheta$-Funktionen:

\weq{7}{
\vartheta_{00}(v),\quad \vartheta_{01}(v),\quad \vartheta_{10}(v),\quad \vartheta_{11}(v),
}

von denen die erste auch mit $\vartheta(v)$ bezeichnet wird.

\gesperrt{In der Folge werden unter Charakteristiken und
$\vartheta$-Funktionen nur noch Hauptcharakteristiken und
Hauptfunktionen verstanden.}

